% ============================================================
% sec/05analisis_tecnico.tex
%
% Análisis exploratorio, diccionario de datos, calidad e
% imputación justificadas, valores atípicos, desbalance de
% clases y prevención explícita de fuga de información.
% ============================================================
\section{Análisis Técnico}
\label{sec:analisistecnico}

El enfoque adoptado es el aprendizaje automático clásico sobre datos
tabulares, según la justificación desarrollada en la
Sección~\ref{subsec:ejeb}. Esa elección impone un conjunto específico de
análisis previos al modelado: caracterizar la calidad de los datos, definir
el tratamiento de los valores ausentes y de los extremos, medir el balance
de clases, y demostrar que la fuga de información está controlada. Esta
sección desarrolla esos análisis.

Todas las cifras que siguen proceden de una auditoría del corpus ejecutada
antes de cualquier modelado, bajo una semilla fija y reproducible. Las
magnitudes que aún no han sido medidas se nombran como pendientes y no se
estiman.

\subsection{Análisis exploratorio}
\label{subsec:eda}

\subsubsection{Alcance de lo medido}

El corpus completo abarca aproximadamente dos meses de publicaciones
diarias, con un volumen del orden de cientos de miles de partidas y más de
un terabyte sin comprimir. La auditoría cubre una muestra temporal de diez
días, elegidos para abarcar las fases inicial, media y tardía de la
competencia, con tres días consecutivos al comienzo de modo que pudiera
medirse la variación entre días contiguos.

Dentro de esa muestra se analizaron 52\,085 partidas a partir de unos
180~GiB de datos, sin errores de lectura, y se obtuvieron 7\,038\,378
observaciones utilizables. Toda afirmación de esta sección está acotada a esa
muestra y no caracteriza los días no auditados.

\subsubsection{Estructura y unidad de observación}

La duración de las partidas está fuertemente sesgada a la derecha. El número
de pasos por partida tiene mediana 140, percentil 95 de 221 y máximo 6\,900.
El recuento de turnos es mucho menor ---mediana 13, percentil 95 de 27,
máximo 1\,035--- porque un solo turno se descompone en muchas acciones
atómicas.

Esa brecha determina cómo debe definirse una etapa de la partida. Fijar el
juego temprano en los turnos uno a ocho sería arbitrario contra una
distribución cuya mediana es 13, de modo que los cortes se derivan de la
distribución empírica y se declaran antes de examinar el conjunto de prueba.

La unidad de observación es el estado del jugador que actúa, y la decisión es
medida antes que estilística. En un sondeo de 139\,672 estados, la
observación correspondiente al jugador que no actúa resultó idéntica a su
propio estado anterior en el 98.8\,\% de los casos. Tratar ambos lados de
cada paso como filas separadas duplicaría casi todas las observaciones sin
añadir información, y repesaría las partidas por su duración.

Siete millones de observaciones frente a cincuenta y dos mil partidas es el
hecho estructural del que depende todo lo que sigue.

\subsubsection{Distribución del desenlace}

El desenlace se resuelve en el 99.86\,\% de las partidas auditadas. A nivel
de partida, el primer jugador vence en el 52.10\,\% de los casos y el
segundo en el 47.70\,\%, con un 0.063\,\% de empates y un 0.140\,\% de
partidas abandonadas. Restringido a las partidas con ganador decidido, la
tasa de victoria del jugador que mueve primero es 0.5287.

\subsubsection{Variación entre días}

Los diez días auditados no son intercambiables. Medida con la divergencia de
Jensen-Shannon entre el primer y el último día, la distribución de mazos es
esencialmente disjunta ---1.0000--- y la de agentes casi tanto ---0.9723---,
mientras que la duración de las partidas apenas varía ---0.1221.

La diversidad no decae de forma suave: los mazos distintos por día van de
1\,268 a 48 entre dos días concretos de la muestra, y los agentes distintos
de 2\,309 a 79 en esos mismos días. La mediana de la puntuación de los
agentes seleccionados sube de 628 el primer día a 1\,141 en un día tardío.

Varias versiones del motor coexisten en la muestra y están ordenadas en el
tiempo, lo que significa que el tiempo de calendario y la versión del motor
están confundidos. La auditoría no puede separar la evolución de las
estrategias del cambio del simulador, y este trabajo reporta esa confusión en
lugar de resolverla. Tres corpus comparables controlaron esa variable de
forma explícita
\cite{hodge_winprediction_2021,xenopoulos_winprob_2022,saravanan_metadiscovery_2024};
no poder hacerlo aquí es una limitación propia.

Una consecuencia de esa variación amenaza a una característica concreta: la
probabilidad de que el jugador que mueve primero ocupe la primera posición
en la mesa era 0.5089 en un día temprano y 0.9991 desde un día intermedio en
adelante. En los días tardíos, posición y orden de turno son casi la misma
variable, y de ahí la normalización de perspectiva descrita en la
Sección~\ref{subsec:pipeline}.

\subsection{Diccionario de datos}
\label{subsec:diccionario}

Un censo sobre 139\,672 estados encontró 34 campos presentes en la totalidad
de las observaciones inspeccionadas: 11 globales, 13 por jugador y 10 por
Pok\'emon. El diccionario es preliminar y se revisará tras el análisis
exploratorio de la tabla materializada.

\begin{table*}[!t]
\centering
\caption{Diccionario de datos por familia de características. La presencia
es la proporción medida de observaciones inspeccionadas en las que el campo
de origen existe.}
\label{tab:diccionario}
\begin{threeparttable}
\footnotesize
\setlength{\tabcolsep}{4pt}
\begin{tabularx}{\textwidth}{@{}
    P{2.2cm}
    >{\centering\arraybackslash}p{1.9cm}
    W{0.80}
    W{1.20}
  @{}}
\toprule
\textbf{Familia} & \textbf{Tipo} & \textbf{Campos y derivación} &
\textbf{Presencia y riesgo}\\
\midrule

Progreso & entero, binario &
Número de turno y contador de acciones, directos; indicador de jugador
inicial; indicadores de energía asignada y de retirada &
Presencia 100\,\%. El indicador de jugador inicial requiere atención:
posición y orden de turno están casi perfectamente confundidos en los días
tardíos
\\
\addlinespace[3pt]

Recursos & entero &
Cartas restantes en mazo y en mano, directas; tamaño del descarte, derivado &
Presencia 100\,\%
\\
\addlinespace[3pt]

Premios & entero &
Cartas de premio restantes, derivadas contando las ranuras no vacías &
Presencia 100\,\%; el contenido está siempre oculto
\\
\addlinespace[3pt]

Tablero & entero, continuo, binario &
Ocupación de la banca; puntos de vida actuales, máximos y su razón; totales
visibles; Pok\'emon dañados; energías, cartas de energía y herramientas;
indicadores de evolución y de aparición en el turno &
Presencia 100\,\%. La razón de puntos de vida debe respetar
$0 \leq hp \leq hp_{\max}$, y esa verificación forma parte del control de
calidad
\\
\addlinespace[3pt]

Estados alterados & binario &
Envenenado, quemado, dormido, paralizado y confundido, por jugador; estadio
en juego y cartas de apoyo ya jugadas, globales &
Presencia 100\,\%
\\
\addlinespace[3pt]

Diferenciales & entero, continuo &
Valor propio menos valor del oponente para cartas en mano, premios, banca,
puntos de vida, energía y mazo &
Derivados. Construcción análoga a la empleada por Hodge \textit{et al.}
\cite{hodge_winprediction_2021}
\\
\addlinespace[3pt]

Contexto de mazo & categórico &
Identificador del mazo sobre un catálogo de 2\,376 mazos distintos &
Cardinalidad elevada y riesgo alto de memorizar el conjunto de estrategias
dominante. Se admite únicamente en la ablación declarada, con mazos
reservados
\\
\midrule

Agrupación \newline \textit{(nunca características)} & --- &
Identificador de partida, fecha y número de turno &
Se conservan para agrupar, estratificar y reportar por etapa; excluidos del
vector predictivo por construcción
\\

\bottomrule
\end{tabularx}
\begin{tablenotes}[flushleft]
\footnotesize
\item \textit{Nota.} Dos campos se excluyen sobre bases medidas antes de
cualquier modelado. El campo de resultado del estado actual es constante en
la totalidad de las observaciones inspeccionadas y no transporta información
en el momento de la predicción. El campo de observación en curso es nulo en
el 98.1\,\% de ellas, y esa nulidad es una propiedad del juego y no un
defecto, según argumenta la Sección~\ref{subsec:imputacion}.
\end{tablenotes}
\end{threeparttable}
\end{table*}

Hodge \textit{et al.} incorporan además, para cada característica, el cambio
respecto del instante anterior \cite{hodge_winprediction_2021}. Este trabajo
no adopta esa construcción. La pregunta planteada concierne a la información
contenida en un estado aislado, y añadir variación temporal introduciría un
vector de fuga adicional ---verificar que el instante anterior sea
estrictamente anterior--- cuya comprobación excede el alcance de esta etapa.
Queda declarada como extensión natural.

\subsection{Calidad de los datos}
\label{subsec:calidad}

\subsubsection{Verificaciones realizadas}

Las 52\,085 partidas auditadas se leyeron sin error. Una única firma de
esquema cubre la totalidad de ellas, de modo que el corpus es
estructuralmente homogéneo pese a las varias versiones del motor. Los
índices de turno resultaron monótonos en todas las partidas inspeccionadas.

Tres comprobaciones independientes de duplicidad devolvieron cero: ningún
identificador de partida repetido, ningún resumen criptográfico repetido
sobre los ficheros, y ninguna huella canónica repetida, calculada esta
última sobre la semilla de configuración, ambos mazos en orden, la secuencia
completa de acciones y el desenlace. Tampoco se registró ninguna partida de
un agente contra sí mismo.

La duplicación exacta no es, por tanto, una preocupación. La duplicación
funcional dentro de una partida es otra cuestión: las observaciones
consecutivas son casi idénticas por construcción. Eso no es un defecto que
deba eliminarse ---es la estructura de dependencia de los datos--- y se
aborda agrupando, no deduplicando.

\subsubsection{Verificaciones pendientes}

La auditoría verificó el corpus crudo. Sobre la tabla materializada deben
reverificarse, y se consignan aquí como pendientes antes que como supuestos:
la ausencia de valores por columna; los rangos imposibles ---puntos de vida
negativos o superiores al máximo, recuentos negativos, ocupación de banca
por encima de su límite---; la consistencia lógica entre campos
relacionados; las filas duplicadas; y la cardinalidad efectiva de las
columnas categóricas.

\subsection{Valores ausentes e imputación}
\label{subsec:imputacion}

La política distingue dos mecanismos, porque tratarlos por igual sería un
error. Emmanuel \textit{et al.} muestran que la eficacia relativa de los
métodos de imputación se invierte según el conjunto evaluado
\cite{emmanuel_missing_2021}: no existe una técnica universalmente superior,
y la estrategia debe derivarse del mecanismo de ausencia y no de la
costumbre.

\noindent\textbf{Ausencia estructural.} Un campo es nulo porque las reglas
del juego determinan que no exista un valor en ese estado: no hay Pok\'emon
activo, no hay estadio en juego, o la carta está oculta al jugador. El caso
más frecuente es el campo de observación en curso, nulo en el 98.1\,\% de
las observaciones inspeccionadas, y el segundo es la mano del oponente, nula
en la totalidad de ellas.

Estos casos no son valores estadísticamente ausentes y no se imputan. Se
codifican con un indicador binario de presencia más un valor centinela
explícito. Imputar una media aquí fabricaría un estado de juego que no puede
ocurrir, y destruiría además la señal que el propio campo transporta: que no
haya nada que observar es información sobre la situación.

\noindent\textbf{Ausencia inesperada.} Un campo requerido falta porque una
partida está malformada o porque el procesamiento es incompleto. El
procedimiento mide primero la frecuencia. Si resulta residual, se excluye la
observación o la partida dejando constancia del criterio y del recuento. Si
resulta sistemática, se trata como defecto del procesamiento o cambio de
esquema y se corrige la tubería, en lugar de parchear los datos.

Ninguna variable numérica se imputa de forma automática antes de haber
caracterizado su mecanismo de ausencia. Sea cual sea el mecanismo, el
estimador de imputación se ajusta dentro del pliegue de entrenamiento: la
imputación figura entre los pasos de preprocesamiento que producen fuga
cuando se estiman sobre el conjunto completo antes de particionar
\cite{sasse_leakage_2025}.

El modelo B4 seleccionado maneja además los valores ausentes de forma nativa
\cite{borisov_tabular_2024}, lo que reduce la dependencia del resultado
respecto de la estrategia adoptada.

\subsection{Valores atípicos}
\label{subsec:atipicos}

En este corpus un valor extremo suele ser un atípico válido antes que un
error. El recuento de turnos alcanza 1\,035 frente a una mediana de 13, y
una partida de esa duración merece investigación y no borrado: puede
corresponder a un tiempo agotado, a un ciclo de acciones, a una interacción
inusual de cartas, a un fallo del simulador o a una partida genuinamente
larga, y esos cinco casos exigen respuestas distintas.

El procedimiento es diagnóstico antes que correctivo. Los candidatos se
marcan por dos vías separadas, que significan cosas distintas: distancia
distribucional sobre las características de duración y recuento, y rangos
imposibles según las reglas del juego. Las partidas marcadas se inspeccionan
caso por caso, y solo se eliminan los registros demostrablemente corruptos.

Cualquier regla de exclusión se declara antes de examinar el conjunto de
prueba. Una regla elegida después de observar su efecto sobre la puntuación
de prueba constituiría fuga por decisión de diseño
\cite{kaufman_leakage_2011}.

Cabe declarar un desacuerdo con la práctica de un corpus comparable: los
autores de SC2EGSet reemplazaron por la mediana los valores que excedían
tres desviaciones estándar \cite{bialecki_sc2egset_2023}. No se adopta esa
regla aquí, porque en un juego por turnos la cola de la distribución de
duración puede ser exactamente donde reside el comportamiento interesante.

\subsection{Desbalance de clases}
\label{subsec:desbalance}

A nivel de partida el desenlace está cerca del equilibrio: 52.10\,\% frente
a 47.70\,\%. Esa no es, sin embargo, la distribución que verá el modelo. La
etiqueta se define desde la perspectiva del jugador que actúa, y cada
partida aporta un número distinto de observaciones desde cada lado, de modo
que la prevalencia sobre la tabla final es función de la tabularización y
aún no ha sido medida. Debe recalcularse sobre la tabla materializada antes
de tomar decisión alguna; presentar la proporción por partida como si fuera
la proporción por fila reportaría como medido algo que no lo está.

La regla de decisión queda fijada de antemano, de modo que no pueda elegirse
para acomodar un resultado.

\noindent\textbf{Si el desbalance resulta leve, no se corrige.} Corregir un
desbalance que no lo necesita distorsiona las probabilidades predichas, y
esas probabilidades son la salida de este trabajo y no un subproducto. Van
den Goorbergh \textit{et al.} muestran que el submuestreo, el sobremuestreo
y la generación sintética no mejoran de forma sistemática la capacidad de
ordenar, y sí sobreestiman de manera pronunciada la probabilidad de la clase
minoritaria \cite{vandenGoorbergh_imbalance_2022}. Conviene precisar el
alcance de esa referencia: no evalúa la ponderación de clases, de modo que
no puede invocarse para preferir esa técnica.

\noindent\textbf{Si resulta apreciable, se comparan alternativas} dentro del
pliegue de entrenamiento, y se reporta su efecto tanto sobre la capacidad de
ordenar como sobre la calidad de las probabilidades, dado que el desbalance
afecta a ambas \cite{silvafilho_calibration_2023}. No se emplea
sobremuestreo sintético sin justificación explícita, y en ningún caso fuera
del pliegue de entrenamiento \cite{sasse_leakage_2025}.

\subsection{Prevención de la fuga de información}
\label{subsec:leakage}

De esta subsección depende la validez de todo lo demás.

\subsubsection{La estructura de dependencia, medida}

Cada partida aporta en promedio unas 135 observaciones, y dos observaciones
consecutivas difieren en una sola acción atómica. Las filas son por tanto
fuertemente dependientes dentro de una partida e independientes solo entre
partidas.

La consecuencia se midió directamente: bajo una partición aleatoria por
observación, la fracción esperada de partidas presentes en ambos
subconjuntos es 0.999945. El protocolo ingenuo contaminaría, en la práctica,
la totalidad del corpus.

\subsubsection{Qué campos pueden usarse}

Un campo puede existir en los datos y aun así ser inválido como
característica. El criterio que gobierna es que el modelo solo puede emplear
información de un instante no posterior al de la predicción
\cite{kaufman_leakage_2011}, afinado por la pregunta que Sasse \textit{et
al.} recomiendan hacerle a cada columna: si ese campo estaría disponible en
el momento real de la decisión \cite{sasse_leakage_2025}. Aplicar esa prueba
campo por campo produce cuatro clases, representadas en la
Figura~\ref{fig:clasescampos}.

\begin{figure}[!t]
\centering
\begin{tikzpicture}[
  font=\scriptsize,
  ok/.style={draw=tecsteel, fill=boxbg, rounded corners=2pt,
             text width=3.3cm, align=center, inner sep=4pt,
             minimum height=7mm},
  duda/.style={draw=decisionambar, fill=decisionambar!8, rounded corners=2pt,
               text width=3.3cm, align=center, inner sep=4pt,
               minimum height=7mm},
  mal/.style={draw=red!60!black, fill=red!5, rounded corners=2pt,
              text width=3.3cm, align=center, inner sep=4pt,
              minimum height=7mm},
  det/.style={font=\tiny\itshape, text=labelgray, align=left,
              text width=3.9cm},
  eje/.style={font=\tiny\bfseries, text=tecnavy}
]

\node[eje] at (-2.05,2.1) {CLASE};
\node[eje] at ( 1.85,2.1) {EJEMPLOS Y DESTINO};

\node[ok]   (a) at (-2.05,1.55) {Seguras};
\node[det]  at ( 1.85,1.55) {estado propio en el turno $\tau$ y sus
                             derivados $\rightarrow$ entran};

\node[duda] (b) at (-2.05,0.75) {Requieren investigación};
\node[det]  at ( 1.85,0.75) {indicador de jugador inicial; registro de
                             acciones $\rightarrow$ solo si se trunca
                             antes de $\tau$};

\node[mal]  (c) at (-2.05,-0.15) {Fuga directa};
\node[det]  at ( 1.85,-0.15) {desenlace; mazo restante ordenado; marcas de
                              tiempo; nombres de agente $\rightarrow$
                              excluidas};

\node[mal]  (d) at (-2.05,-1.05) {Constantes};
\node[det]  at ( 1.85,-1.05) {campo de resultado del estado actual, fijo en
                              todas las observaciones $\rightarrow$
                              excluida};

\draw[draw=tecrule, dashed] (-0.05,2.0) -- (-0.05,-1.5);

\end{tikzpicture}
\caption{Clasificación de los campos según la prueba de disponibilidad en el
momento de la decisión. Solo la primera clase alimenta el vector de
características; la segunda requiere verificación adicional antes de
admitirse.}
\label{fig:clasescampos}
\end{figure}

Dos exclusiones merecen justificación explícita. El campo que expone el mazo
restante de cada jugador en orden está presente en la totalidad de las
partidas y constituye información perfecta que ningún jugador posee en el
momento de decidir. Y los nombres de agente actúan como sustituto de la
habilidad: la mediana de puntuación de los agentes seleccionados sube de 628
a más de 1\,100 a lo largo de la muestra, de modo que un identificador de
agente codifica parcialmente el desenlace.

\subsubsection{La información oculta ya está protegida}

El entorno retiene lo que debe retener, y se comprobó en lugar de suponerlo.
Sobre 139\,672 observaciones inspeccionadas, la mano del oponente fue nula
en la totalidad de los casos, mientras que su recuento de cartas estuvo
siempre disponible como entero. Las 1\,281\,678 ranuras de premio
inspeccionadas ---propias y del oponente por igual--- fueron todas nulas.

La fuga de información oculta no es, por tanto, un riesgo que deba
mitigarse: es una propiedad que no debe destruirse. De ahí que las
características se extraigan exclusivamente de la observación propia del
jugador activo.

\subsubsection{Vectores de riesgo y contramedidas}

\begin{table}[!t]
\centering
\caption{Vectores de fuga, evidencia de auditoría que establece cada uno
como riesgo real, y control aplicado.}
\label{tab:leakage}
\footnotesize
\begin{tabularx}{\columnwidth}{@{}P{1.5cm} P{2.3cm} Y@{}}
\toprule
\textbf{Vector} & \textbf{Evidencia medida} & \textbf{Control}\\
\midrule
Por partida &
Una partición aleatoria por fila sitúa 0.999945 de las partidas en ambos
subconjuntos &
Partición y pliegues agrupados por partida; prueba de intersección
vacía\\[2pt]
Información futura &
El mazo restante ordenado está presente en el 100\,\% de las partidas &
Características extraídas solo del estado propio en $\tau$; campo
excluido\\[2pt]
Sustituto de habilidad &
La mediana de puntuación de los agentes sube de 628 a más de 1\,100 &
Metadatos de agente usados solo para auditoría y agrupación\\[2pt]
Memorización de mazos &
2\,376 mazos, pero los diez más frecuentes cubren el 39.94\,\%; un día
contiene solo 48 &
Escenario C; contexto de mazo confinado a una ablación declarada\\[2pt]
Información oculta &
Cero fugas en 139\,672 observaciones; 1\,281\,678 ranuras de premio todas
nulas &
Lectura restringida a la observación propia, con verificación
automática\\[2pt]
Preprocesamiento &
Riesgo de diseño, no propiedad del corpus &
Transformaciones ajustadas dentro del pliegue de entrenamiento
\cite{sasse_leakage_2025}\\[2pt]
Temporal &
Divergencia de mazos de 1.0000 entre el primer y el último día auditado &
El escenario B corta por fecha, nunca al azar
\cite{kaufman_leakage_2011}\\
\bottomrule
\end{tabularx}
\end{table}

Acompañar cada vector con evidencia medida sobre este corpus, y no solo con
la descripción del riesgo, es lo que distingue una prevención verificada de
una precaución declarada.

\subsubsection{Verificación automática}

La tubería incorpora invariantes que deben fallar antes de cualquier
entrenamiento. Primero, ninguna columna excluida puede aparecer en el
conjunto de características. Segundo, la intersección de identificadores de
partida entre cualquier par de subconjuntos debe ser vacía. Tercero, toda
carta del corpus debe estar presente en el conjunto de entrenamiento, según
el procedimiento de la Sección~\ref{subsec:protocolo}. Y cuarto, ninguna
transformación puede haberse ajustado con datos de validación, calibración o
prueba.

Una política de prevención sin una prueba que la verifique constituye una
intención y no una mitigación.

\subsubsection{Medición de la inflación atribuible a la fuga}

Se ejecutará el modelo B2 bajo dos particiones ---aleatoria por fila y
agrupada por partida--- y se reportará la diferencia. Ese procedimiento
convierte la prevención de fuga en un resultado medido sobre los propios
datos, siguiendo el diseño con que Figueiredo y Mendes cuantifican el efecto
en su dominio \cite{figueiredo_leakage_2024}.

Se emplea únicamente B2 porque el punto que se demuestra no depende del
modelo, y porque Karbalaie \textit{et al.} encuentran que los modelos
lineales son los menos inflados por un protocolo defectuoso
\cite{karbalaie_participantaware_2026}: si B2 muestra inflación, el valor
obtenido es una cota inferior del efecto.

La partición aleatoria se emplea exclusivamente como demostración del sesgo.
Ningún resultado del trabajo se reporta bajo ella.

\subsection{Representatividad}
\label{subsec:representatividad}

El corpus no es una muestra aleatoria del juego. Las repeticiones se
seleccionan a diario ordenadas por puntuación de los agentes y se truncan a
un volumen fijo \cite{ptcg_dataset_2026}, y de ahí se siguen tres
consecuencias, las tres medidas.

La población es de agentes automáticos y no de personas. Corresponde a la
franja alta de la clasificación, con la puntuación mediana ascendiendo a lo
largo de la muestra. Y está concentrada: los diez mazos más frecuentes suman
el 39.94\,\% de las instancias y los cien más frecuentes el 74.76\,\%, con
días individuales que descienden a 48 mazos distintos.

Los dos corpus análogos acotan este caso desde direcciones opuestas. SC2EGSet
contiene únicamente repeticiones de nivel de torneo y sus autores declaran
que no admite análisis entre niveles de habilidad
\cite{bialecki_sc2egset_2023}; el certamen sobre \emph{Hearthstone} generó
sus estados con agentes que actúan al azar
\cite{janusz_hearthstone_2017}. Uno procede de la población más fuerte
posible y el otro de la más débil, y ninguno es la población general.
Cualquier afirmación de este trabajo es, en consecuencia, una afirmación
sobre agentes automáticos competitivos de una franja de puntuación
específica durante una muestra específica de diez días.